\scanpage{91}
\mspage{42}

\deleted{Une \unsure{filtration|stratification} \inserted{\deleted{\unsure{ordinaire|indicielle}}} de type $I$}
Soient $X$ un espace, $I$ un ens. ordonné.
Une \uline{stratification \unsure{ordinaire|indicielle}}, de type $I$, de $X$, est une application
\[I\longrightarrow\mathcal P(X)\qquad i\longmapsto X_i\]
satisfaisant les deux conditions

(Strat 1) les $X_i$ fermés, la famille des $X_i$ \inserted{est un \uline{recouvrement}} \uline{\unsure{loc\textsuperscript{t}|loc!} \unsure{fini|finie}}\inserted{ de $X$.}

(Strat 2) $i\leq j\Longrightarrow X_i\subset X_j$ i.e. $i\mapsto X_i$ est croissante

(Strat 3) $\forall i,j\in I$, on a
\[X_i\cap X_j=\bigcup_{k\in\bar i\cap\bar j}X_k.\]
\margin{les $X_i$ sont les ``strates'' \uline{fermées}}

\deleted{\unsure{Pour ces conditions|Posons les conditions}}
Pour tout drapeau
\[d=(i_0<\cdots<i_n)\]
de $I$, posons
\[(88)\qquad X_d=X_{i_0},\]
alors \deleted{\unsure{$\varinjlim$|les}} $X_d{}_{d\in\mathrm{Drap}I}$ forment un \deleted{\unsure{syst.|objet}} espace semi-simplicial \deleted{\unsure{large}} au sens large (pour des applications strictement croissantes sur les simplexes types)\inserted{, au dessus de $X$}, qui fournit en un sens convenable une ``résolution'' semi-simpl. de $X$. On aura p.ex. pour tt faisceau abélien $F$ sur $X$
\[
(89)\qquad\begin{aligned}
H^*(X,F)\Longleftarrow E_2^{pq}&=H^p(I,i\mapsto H^q(X_i,F))\\
&=H^p_{ss}(\mathrm{Drap}(I),d\mapsto H^q(X_d,F))
\end{aligned}
\]
Mais les \inserted{$X_i$,} $X_d$ sont des espaces souvent compliqués, et le dévissage (89), souvent n'est pas pratique, on remplace par les $X_i^*$, $X_d^*=X_{i_0}^*$, en posant
\margin{90}
